\section*{Capítulo \ref{ch:g11:mechanical-energy} --- La energía mecánica y su conservación}

\begin{solution}{exo:g11:mechanical-energy:1}
(a) $v = \qty{6.94}{m/s}$:
$E_k = \tfrac12 \times 90 \times 6.94^2 \approx \qty{2.2e3}{J}$.
(b) $v = \qty{36.1}{m/s}$:
$E_k = \tfrac12 \times 1200 \times 36.1^2 \approx \qty{7.8e5}{J}$.
(c) $\times 4$; $\times 9$ --- el cuadrado.
\end{solution}

\begin{solution}{exo:g11:mechanical-energy:2}
$\Delta E_p = 75 \times 9.81 \times 850 \approx \qty{6.3e5}{J}$. No: la
referencia desplaza por igual los dos valores. Tampoco: solo entra la
diferencia de altura.
\end{solution}

\begin{solution}{exo:g11:mechanical-energy:3}
$v = \sqrt{2 \times 9.81 \times 20} = \qty{19.8}{m/s} \approx
\qty{71}{km/h}$. En $mgh = \tfrac12 mv^2$ la masa se cancela.
\end{solution}

\begin{solution}{exo:g11:mechanical-energy:4}
$W = \Delta E_k = 0 - \tfrac12 \times 1200 \times 25^2 =
\qty{-3.75e5}{J}$; $F = \num{3.75e5}/50 = \qty{7.5e3}{N}$.
\end{solution}

\begin{solution}{exo:g11:mechanical-energy:5}
$v = \sqrt{2 \times 9.81 \times 0.080} \approx \qty{1.3}{m/s}$; vuelve a
subir hasta \qty{8.0}{cm} (conservación).
\end{solution}

\begin{solution}{exo:g11:mechanical-energy:6}
$\num{1.5e9}/\num{2.6e3} \approx \num{5.9e5}$: unas seiscientas mil
balas. $h = 88.9^2/(2 \times 9.81) \approx \qty{4.0e2}{m}$ --- su propia
rapidez podría elevar al tren $400$ \omterm{def:g10:orders-of-magnitude:unit}{metros}.
\end{solution}

\begin{solution}{exo:g11:mechanical-energy:7}
La misma rapidez,
$v = \sqrt{2 \times 9.81 \times 3.0} \approx \qty{7.7}{m/s}$ (solo entra
el desnivel). El tobogán empinado gana en tiempo: alcanza antes una
rapidez alta y su camino es más corto.
\end{solution}

\begin{solution}{exo:g11:mechanical-energy:8}
$h = v^2/2g = 144/19.62 \approx \qty{7.3}{m}$. A \qty{4.0}{m}:
$v = \sqrt{144 - 2 \times 9.81 \times 4.0} \approx \qty{8.1}{m/s}$ ---
la misma en los dos sentidos: $E_m$ depende de la altura, no del sentido
de la marcha.
\end{solution}

\begin{solution}{exo:g11:mechanical-energy:9}
$v = \sqrt{23^2 + 2 \times 9.81 \times 35} = \sqrt{1216} \approx
\qty{35}{m/s}$. En la realidad, menor: la resistencia del aire disipa
parte de $E_m$.
\end{solution}

\begin{solution}{exo:g11:mechanical-energy:10}
$\Delta E_m = \tfrac12 \times 45 \times 8.0^2 - 45 \times 9.81 \times
6.0 = 1440 - 2649 \approx \qty{-1.2e3}{J}$;
$f = 1209/15 \approx \qty{81}{N}$. Ha ido a calor, en los patines y en
la nieve.
\end{solution}

\begin{solution}{exo:g11:mechanical-energy:11}
$v = \sqrt{2 \times 9.81 \times 0.15} \approx \qty{1.7}{m/s}$; vuelve a
subir hasta \qty{15}{cm}. El clavo cambia el camino (una circunferencia
más cerrada), no la \omterm{def:g10:energy-conservation:energy}{energía}: las alturas y las rapideces no se tocan.
\end{solution}

\begin{solution}{exo:g11:mechanical-energy:12}
$v = \sqrt{2 \times 9.81 \times 2.5} = \qty{7.0}{m/s}$. Las alturas de
subida escalan con $E_m$, así que quedan multiplicadas por $0.9$ en cada
travesía: $2.25$, $2.03$, $1.82$, $1.64$, $1.48$, $1.33$, $1.20$,
$1.08$, $0.97$ --- la primera por debajo de \qty{1.0}{m} es la $9$.ª
travesía.
\end{solution}

\begin{solution}{exo:g11:mechanical-energy:13}
La $E_m$ inicial vale
$\tfrac12 m \times 15^2 + m \times 9.81 \times 45$ sea cual sea el
ángulo, así que la rapidez al llegar al suelo es
$v = \sqrt{225 + 2 \times 9.81 \times 45} = \sqrt{1108} \approx
\qty{33}{m/s}$. El ángulo fija la \emph{dirección} de la velocidad de
llegada, el tiempo de vuelo y el \omterm{def:g11:fundamental-interactions:strong}{alcance} --- no la rapidez de llegada.
\end{solution}

\begin{solution}{exo:g11:mechanical-energy:14}
$\tfrac12 mv^2 = mgh - f\ell = 80 \times 9.81 \times 120 - 65 \times
800 = \num{94176} - \num{52000} = \qty{4.22e4}{J}$, luego
$v = \sqrt{2 \times \num{42176}/80} \approx \qty{32}{m/s}$
(\qty{117}{km/h}). Fracción disipada:
$52000/94176 \approx 55\,\%$.
\end{solution}

\begin{solution}{exo:g11:mechanical-energy:15}
$m = \qty{2.0e9}{kg}$:
$E_p = \num{2.0e9} \times 9.81 \times 300 \approx \qty{5.9e12}{J}
= \num{5.9e12}/\num{3.6e6} \approx \qty{1.6e6}{kW.h}$. Recuperados:
$0.85 \times \num{1.64e6} \approx \qty{1.4e6}{kW.h}$ --- unos
$\num{1.4e5}$ hogares durante un día.
\end{solution}

\begin{solution}{pb:g11:mechanical-energy:1}
\textbf{1.} $v_{\text{top}} = \sqrt{gR} = \sqrt{9.81 \times 8.0}
\approx \qty{8.9}{m/s}$.

\textbf{2.} $E_m = mg\,(2R) + \tfrac12 m gR = \num{3.14e5} +
\num{7.8e4} \approx \qty{3.92e5}{J}$.

\textbf{3.} De $mgH_{\min} = E_m$ sale $H_{\min} = \qty{20}{m}$; en
álgebra, $H_{\min} = 2R + \frac{gR}{2g} = \frac{5R}{2}$.

\textbf{4.} $v_{\text{top}} = \sqrt{2g(H - 2R)} =
\sqrt{2 \times 9.81 \times 8.0} \approx \qty{12.5}{m/s}$;
$v_{\text{top}}^2/(gR) = 2.0$ --- el doble del mínimo.

\textbf{5.} Todas las \omterm{def:g10:energy-conservation:energy}{energías} son proporcionales a $m$, así que $m$ se
cancela: las mismas rapideces y la misma seguridad, lleno o vacío.

\textbf{6.} $v = \sqrt{2 \times 9.81 \times 24} \approx \qty{21.7}{m/s}
\approx \qty{78}{km/h}$.

\textbf{7.} $v = \sqrt{2 \times 9.81 \times (24 - 8.0)} \approx
\qty{17.7}{m/s}$.

\textbf{8.} $mgH = \qty{4.71e5}{J}$. $(E_p, E_k)$: lanzamiento
$(\num{4.71e5}, 0)$; base $(0, \num{4.71e5})$; costado
$(\num{1.57e5}, \num{3.14e5})$; cima
$(\num{3.14e5}, \num{1.57e5})$ --- cada pareja suma \qty{4.71e5}{J}.

\textbf{9.} $v = \sqrt{2 \times 9.81 \times (24 - 12)} \approx
\qty{15.3}{m/s}$.

\textbf{10.} $v(z) = \sqrt{2g(H - z)}$: lo más rápido, en el punto más
bajo de la vía; lo más lento, en la cima del rizo.

\textbf{11.} A la misma altura, $\Delta E_p = 0$ y
$\Delta E_m = \Delta E_k = \tfrac12 \times 2000 \times (19.2^2 -
20.4^2) \approx \qty{-4.75e4}{J}$: la colocación hace que el déficit sea
puramente cinético, leído en dos velocímetros.

\textbf{12.} $W(\vect f) = \qty{-4.75e4}{J}$;
$f = \num{4.75e4}/60 \approx \qty{7.9e2}{N}$.

\textbf{13.} \omterm{def:g10:energy-conservation:forms}{Energía térmica}, en las ruedas, en los raíles y en el aire.

\textbf{14.} $E_m$ en el primer sensor:
$\num{4.16e5} + \num{3.9e4} = \qty{4.55e5}{J}$; pérdida
$\num{4.75e4}/\num{4.55e5} \approx 10\,\%$ por cada \qty{60}{m}. Con
$H_{\min}$, la cima del rizo se alcanza \emph{justo} a $\sqrt{gR}$;
cualquier pérdida la deja por debajo de la rapidez de seguridad.

\textbf{15.} Perdiendo un $\approx 10\,\%$ cada \qty{60}{m}, se ha ido
cerca del $15\,\%$ de $mgH$ ($\approx \qty{7.1e4}{J}$):
$E_m \approx \qty{4.00e5}{J}$, todavía por encima de los
\qty{3.92e5}{J} necesarios --- a salvo, con solo un $2\,\%$ de sobra. El
margen no era un lujo.

\textbf{16.} $E_k = \tfrac12 \times 2000 \times 18.0^2 =
\qty{3.24e5}{J}$.

\textbf{17.} $F = \num{3.24e5}/45 = \qty{7.2e3}{N}$.

\textbf{18.} $a = 7200/2000 = \qty{3.6}{m/s^2} \approx 0.37\,g$: firme
pero cómodo.

\textbf{19.} \omterm{def:g10:inertia:inventory}{Rozamiento} de la vía:
$mgH - \num{3.24e5} = \num{470880} - \num{324000} = \qty{1.47e5}{J}$.
Auditoría: $\num{1.47e5} + \num{3.24e5} = \num{4.71e5} = mgH$ --- todos
los \omterm{def:g11:work-of-force:work}{julios} justificados.

\textbf{20.} La altura de lanzamiento $H = \qty{24}{m}$ fijó la
seguridad del rizo, todas las rapideces $\sqrt{2g(H - z)}$, la factura
del \omterm{def:g10:inertia:inventory}{rozamiento} y la \omterm{def:g10:inertia:force}{fuerza} de los frenos; el precio de ignorar el
\omterm{def:g10:inertia:inventory}{rozamiento} es el margen --- los $\qty{4}{m}$ de colina que el mundo real
consume en silencio.
\end{solution}
